\chapter{Introducción}

\section{Problema y propósito del proyecto}

El agotamiento de procesos, CPU, memoria o almacenamiento puede degradar un
sistema Linux hasta impedir su operación y recuperación. Este riesgo puede
originarse mediante proliferación recursiva de procesos, incluso desde
mecanismos legítimos como el aislamiento por sitio de un navegador
\cite{gierlings2023isolated}. El aislamiento de recursos, por su parte, puede
preservar la interactividad frente a cargas extremas de procesos y memoria
\cite{yang2008redline}.

En este proyecto, \emph{bunny virus} designa un comportamiento sostenido de
agotamiento: creación acelerada de procesos, aumento de memoria o crecimiento
continuo de un archivo. Los cambios de tamaño y otros metadatos permiten
observar actividad sospechosa en almacenamiento \cite{pennington2003storage},
pero una monitorización frecuente también introduce un coste dependiente de la
carga \cite{patil2004i3fs}. Por ello, \emph{anti-bunny-virus} obtiene muestras
de procesos y recursos desde \texttt{/proc}, vigila un directorio configurado y
evalúa tasas temporales y persistencia. Cada alerta conserva su evidencia y la
contención solo puede aplicarse, en laboratorio, a un grupo atribuible y
autorizado.

La pregunta de investigación es: ¿un monitor en C para Linux basado en tasas
temporales, persistencia y correlación puede detectar tempranamente escenarios
controlados de agotamiento y contener de forma segura el grupo responsable? La
evaluación se limita a los simuladores, configuración y máquinas virtuales
descritos; no pretende demostrar protección general contra malware.

\section{Hipótesis}

Se plantea que las tasas temporales y la persistencia permitirán detectar los
tres simuladores antes de sus límites configurados. Asimismo, la validación de
identidad y autorización debería permitir contener únicamente el grupo de
laboratorio, sin alertas críticas durante la línea base. La hipótesis se
evaluará mediante latencia de alerta y contención, recursos consumidos, falsos
positivos, recuperación y sobrecarga del monitor, y sus conclusiones se
restringirán al entorno experimental.

\clearpage

\section{Objetivos}

\subsection{Objetivo general}

Diseñar, implementar y evaluar un monitor en C para Linux capaz de detectar
comportamientos controlados de agotamiento de procesos, memoria y almacenamiento,
y de aplicar una respuesta trazable y segura sobre un grupo autorizado en un
entorno de laboratorio.

\subsection{Objetivos específicos}

\begin{enumerate}
  \item Recolectar desde Linux la relación entre procesos, el consumo de CPU y
  memoria, y el crecimiento de archivos, incluyendo la atribución del archivo a
  un proceso cuando exista evidencia suficiente.

  \item Detectar anomalías sostenidas mediante tasas temporales, umbrales y
  persistencia, y registrar para cada alerta las métricas que originaron la
  decisión.

  \item Implementar una política de contención de laboratorio que valide la
  identidad y autorización del proceso o grupo antes de enviar señales, y que
  rechace acciones sobre objetivos no autorizados.

  \item Comparar experimentalmente los escenarios sin protección, con detección
  y con contención, midiendo latencia de alerta, consumo de recursos, falsos
  positivos, recuperación y sobrecarga del monitor.
\end{enumerate}
